\chapter{正交展开}
\label{chap:orthogonal-expansions}

有限维欧氏空间中，正交基把向量写成坐标之和，并把范数平方写成坐标平方和。
在 Hilbert 空间中，基可能无限，展开因而包含两个问题：系数怎样确定，以及
由这些系数组成的级数是否回到原向量。Bessel 不等式先给出有限部分和的一致
控制；完备正交系再排除“所有系数都为零而向量仍非零”的剩余方向。Parseval
恒等式正是两件事同时成立后的 Pythagoras 等式。

\section{正交系与 Fourier 系数}

设 \(H\) 为 Hilbert 空间。向量族 \((e_j)_{j\in J}\) 称为正交规范系，如果
\(\langle e_j,e_k\rangle=\delta_{jk}\)。对 \(x\in H\)，数
\[
 \widehat x(j)=\langle x,e_j\rangle
\]
称为 \(x\) 关于该系统的 Fourier 系数。名称并不预设系统是三角函数，也不
预设 \(x\) 已经能够展开。

若 \(F\subset J\) 有限，则
\[
 P_Fx=\sum_{j\in F}\langle x,e_j\rangle e_j
\]
是 \(x\) 到 \(\operatorname{span}\{e_j:j\in F\}\) 的正交投影。特别地，
系数不是人为选择的：任何有限线性组合最接近 \(x\) 时，其系数只能是
\(\langle x,e_j\rangle\)。

\section{Bessel 不等式}

由 \(x=P_Fx+(x-P_Fx)\) 的正交分解，
\[
 \|x\|^2=\sum_{j\in F}|\langle x,e_j\rangle|^2+\|x-P_Fx\|^2.
\]
于是得到以下结论。

\begin{theorem}[Bessel 不等式]\label{thm:bessel-inequality}
对任意正交规范系 \((e_j)_{j\in J}\) 和任意 \(x\in H\)，
\[
 \sum_{j\in J}|\langle x,e_j\rangle|^2\le\|x\|^2,
\]
其中对任意指标集的非负和定义为所有有限子和的上确界。
\end{theorem}

这一定义还说明，每个固定 \(x\) 至多有可数多个非零 Fourier 系数。事实上，
对每个正整数 \(m\)，满足
\(|\langle x,e_j\rangle|\ge1/m\) 的指标只能有限；把这些有限集对 \(m\)
取可数并即可。即使正交系不可数，单个向量的有效展开仍是可数级数。

\begin{corollary}\label{cor:orthogonal-series-convergence}
若标量族 \((c_j)_{j\in J}\) 满足
\(\sum_{j\in J}|c_j|^2<\infty\)，则有限和网
\(\sum_{j\in F}c_je_j\) 按有限集包含关系在 \(H\) 中收敛，并且
\[
 \left\|\sum_{j\in J}c_je_j\right\|^2=\sum_{j\in J}|c_j|^2.
\]
\end{corollary}
\begin{proof}
对有限集 \(F,G\)，正交性给
\[
 \left\|\sum_{j\in F}c_je_j-\sum_{j\in G}c_je_j\right\|^2
 =\sum_{j\in F\triangle G}|c_j|^2.
\]
平方和有限意味着右侧在 \(F,G\) 包含某个固定有限集后任意小，故得到
Cauchy 网。只有可数多个系数非零，等价地也可枚举后使用 Hilbert 空间完备性。
范数等式先对有限和成立，再取极限。
\end{proof}

\section{完备正交系}

\begin{definition}
正交规范系 \((e_j)_{j\in J}\) 称为完备的，如果
\[
 \overline{\operatorname{span}}\{e_j:j\in J\}=H.
\]
完备正交系也称 Hilbert 基或正交规范基；这里的“基”允许无限级数，不是
要求每个向量只有有限展开的 Hamel 基。
\end{definition}

\begin{theorem}\label{thm:orthonormal-completeness-equivalences}
对正交规范系 \(E=(e_j)_{j\in J}\)，下列条件等价：
\begin{enumerate}
 \item \(E\) 完备；
 \item 若 \(x\perp e_j\) 对所有 \(j\) 成立，则 \(x=0\)；
 \item 对每个 \(x\in H\)，有限和网 \(P_Fx\) 收敛到 \(x\)；
 \item 对每个 \(x\in H\)，Parseval 恒等式
 \[
   \|x\|^2=\sum_{j\in J}|\langle x,e_j\rangle|^2
 \]
 成立。
\end{enumerate}
\end{theorem}
\begin{proof}
令 \(M=\overline{\operatorname{span}}E\)。正交投影定理给
\(H=M\oplus M^\perp\)，而 \(M^\perp\) 正是与所有 \(e_j\) 正交的向量，
故前两项等价。

对固定 \(x\)，Bessel 不等式说明系数平方可和，因而
\(s=\sum_j\langle x,e_j\rangle e_j\) 存在。对每个 \(k\)，连续取内积得
\(\langle x-s,e_k\rangle=0\)，所以 \(x-s\in M^\perp\)，而
\(s\in M\)。因此 \(s=P_Mx\)。于是 \(M=H\) 当且仅当所有 \(x\) 都有
\(s=x\)，即第一项与第三项等价。最后正交分解给
\[
 \|x\|^2=\|s\|^2+\|x-s\|^2
 =\sum_j|\langle x,e_j\rangle|^2+\|x-s\|^2,
\]
所以第三项与第四项等价。
\end{proof}

“所有系数确定”并不自动表示系统完备。以 \(L^2[-\pi,\pi]\) 为例，若只取
偶函数组成的正交系统，那么每个奇函数与系统中所有元素正交；它们的系数
全为零，但函数本身并非零。必须另证闭线性张成等于整个空间。

\section{Parseval 恒等式与内积展开}

完备时，由极化公式或直接对部分和取极限，还得到
\begin{equation}\label{eq:parseval-inner-product}
 \langle x,y\rangle
 =\sum_{j\in J}\langle x,e_j\rangle
          \overline{\langle y,e_j\rangle}.
\end{equation}
右侧绝对收敛，因为 Cauchy--Schwarz 不等式可先用于有限标量和，再取上确界。
定义分析算子与合成算子
\[
 U:H\to\ell^2(J),\qquad Ux=(\langle x,e_j\rangle)_{j\in J},
 \quad
 S:\ell^2(J)\to H,\qquad S(c)=\sum_jc_je_j.
\]
Bessel 不等式说明 \(\|U\|\le1\)，正交级数定理说明 \(S\) 是到
\(M=\overline{\operatorname{span}}E\) 上的等距映射，并且
\(US=I_{\ell^2(J)}\)、\(SU=P_M\)。所以 \(U\) 总是满射，却只有在系统完备
时才是等距单射；这时 \(S=U^{-1}\)，从而 \(H\) 与 \(\ell^2(J)\) 酉同构。

\section{构造正交基}

对线性无关序列 \(x_1,x_2,\ldots\)，Gram--Schmidt 过程递归定义
\[
 u_n=x_n-\sum_{k<n}\langle x_n,e_k\rangle e_k,
 \qquad e_n=\frac{u_n}{\|u_n\|}.
\]
线性无关保证 \(u_n\ne0\)，且每个前 \(n\) 项的线性张成保持不变。

\begin{theorem}\label{thm:separable-hilbert-basis}
非零 Hilbert 空间可分，当且仅当它有可数完备正交系。
\end{theorem}
\begin{proof}
若 \(e_n\) 是可数完备正交系，以有理实部、虚部系数构成的有限线性组合全体
可数且稠密，故空间可分。

反之，取可数稠密集 \(D=\{x_n\}\)。依次删去落在前面线性张成中的向量，
对剩下的序列作 Gram--Schmidt。所得正交系的闭线性张成包含 \(D\)，因而
等于 \(H\)。若过程有限停止，也得到有限完备正交系。
\end{proof}

一般 Hilbert 空间利用 Zorn 引理可以取极大正交规范系。若其正交补含非零
向量，归一化后还能加入系统，违反极大性；故极大系统必完备。这一论证说明
任意 Hilbert 空间都有正交规范基，但不可分空间的基可能不可数。

\section{广义 Fourier 展开}

设 \(T\in\mathcal L(H)\)，在完备正交系下可用矩阵元
\(t_{jk}=\langle Te_k,e_j\rangle\) 描述。对固定 \(k\)，第 \(k\) 列属于
\(\ell^2(J)\)，且平方和为 \(\|Te_k\|^2\le\|T\|^2\)。然而“每列都在
\(\ell^2\) 且有共同上界”并不足以保证一个无限矩阵定义有界算子；还需控制
不同列叠加后的作用。有限维矩阵直觉在这里第一次出现实质限制。

对闭子空间 \(M\)，若 \((e_j)\) 是 \(M\) 的完备正交系，则
\[
 P_Mx=\sum_j\langle x,e_j\rangle e_j,
 \qquad d(x,M)^2=\|x\|^2-\sum_j|\langle x,e_j\rangle|^2.
\]
这把最佳逼近误差精确写成遗漏系数的能量。三角函数系统的完备性将在
\chapterref{chap:trigonometric-fourier-series}和
\chapterref{chap:l2-fourier-series}借助 Fejér 核证明；在那之前，不能把 Fourier 系数公式本身当作
展开收敛的证明。

\section*{习题}
\addcontentsline{toc}{section}{习题}
\subsection*{A组}
\begin{enumerate}[label=\textbf{\thechapter.\arabic*.}]
\exerciseitem{orthogonal-bessel-finite}从最佳逼近性质直接推出有限正交系的 Bessel 不等式。
\exerciseitem{orthogonal-coeff-countable}证明关于任意正交规范系，一个向量至多有可数多个非零 Fourier 系数。
\exerciseitem{orthogonal-ell-two-synthesis}给定 \(c\in\ell^2\)，证明 \(\sum c_ne_n\) 收敛，并计算任意尾和的范数。
\exerciseitem{orthogonal-complete-kernel}证明正交规范系完备当且仅当没有非零向量与全部基向量正交。
\end{enumerate}
\subsection*{B组}
\begin{enumerate}[label=\textbf{\thechapter.\arabic*.},resume]
\exerciseitem{orthogonal-best-n-term}设 \(e_1,\ldots,e_n\) 正交规范。证明 Fourier 系数给出的组合是所有系数组合中的唯一最佳逼近。
\exerciseitem{orthogonal-parseval-polarization}由 Parseval 恒等式和极化公式证明\eqref{eq:parseval-inner-product}。
\exerciseitem{orthogonal-gram-schmidt-polynomials}在 \(L^2[-1,1]\) 中对 \(1,x,x^2\) 作 Gram--Schmidt 正交化。
\exerciseitem{orthogonal-separable-family}证明可分 Hilbert 空间中每个正交规范集至多可数。
\exerciseitem{orthogonal-projection-tail}若 \(M_n\) 是递增闭子空间且 \(M=\overline{\bigcup_nM_n}\)，证明 \(P_{M_n}x\to P_Mx\)。
\end{enumerate}
\subsection*{C组}
\begin{enumerate}[label=\textbf{\thechapter.\arabic*.},resume]
\exerciseitem{orthogonal-riesz-fischer}证明合成算子 \(S\) 等距映到正交系的闭线性张成，且 \(US=I\)、\(SU=P_M\)。由此刻画 \(U\) 何时为等距同构。
\optionalexerciseitem{orthogonal-haar-first-levels}在 \(L^2[0,1]\) 中验证常函数 \(1\) 与前两层 Haar 函数两两正交规范，并计算 \(f(x)=x\) 的相应系数与投影误差。
\end{enumerate}
